\section{Methodology}
\label{sec:method}

We present the theoretical formulation and mathematical derivation of
\textbf{GP-BayesMerge}. First, we define the parameter-space model-merging
scheme. Second, we introduce the PAC-Bayes generalization framework applied to
test-time merging coefficients. Third, we model the coefficient prior using a
Gaussian Process over network depth and derive the quadratic precision-matrix
regularization. Finally, we formulate the complete multi-task calibration loss.

\subsection{Parameter-Space Model Parameterization}

Let $\theta_0 \in \mathbb{R}^D$ represent the parameters of a pre-trained base
model, and let $\{\theta_k\}_{k=1}^K$ denote the parameters of $K$ distinct
expert models, each fine-tuned from the same base initialization $\theta_0$ on a
specific task or domain. We define the task vector of expert $k$ as $\tau_k =
\theta_k - \theta_0$.

To allow for architectural heterogeneity across the $L$ layers of the network,
we parameterize the merging coefficients as a matrix $\Lambda \in \mathbb{R}^{L
\times K}$, where $\lambda_{l, k} \ge 0$ represents the merging coefficient for
layer $l \in \{1, \dots, L\}$ and task expert $k \in \{1, \dots, K\}$. Let
$\theta_{0, l}$ and $\theta_{k, l}$ be the parameters of layer $l$ of the base
model and expert $k$, respectively. The merged parameters of layer $l$ are given
by:
\begin{equation}
\theta_l(\Lambda) = \theta_{0, l} + \sum_{k=1}^K \lambda_{l, k} (\theta_{k, l} - \theta_{0, l})
\label{eq:param_interp}
\end{equation}
Let $\lambda_k = [\lambda_{1, k}, \dots, \lambda_{L, k}]^T \in \mathbb{R}^L$
denote the spatial coefficient vector for task $k$. Our goal is to optimize the
matrix of coefficients $\Lambda = \{\lambda_k\}_{k=1}^K$ on a small calibration
set during test-time adaptation, yielding the optimized parameters $\Lambda^* = \{\lambda_k^*\}_{k=1}^K$.

\subsection{PAC-Bayes Generalization Bounds over Coefficients}

Instead of placing prior distributions on the high-dimensional weight space
$\mathbb{R}^D$ (which is computationally intractable and ignores architectural
structure), we place a prior distribution directly on the low-dimensional
control space of merging coefficients $\Lambda \in \mathbb{R}^{L \times K}$.

Let $R(\theta(\Lambda))$ be the true target risk of the merged model, and let
$\hat{R}(\theta(\Lambda); X)$ be the empirical risk evaluated on a calibration
batch $X$ of size $N$. Let $P(\Lambda)$ be a prior probability distribution over
the merging coefficients, and let $Q(\Lambda)$ be the optimized posterior
distribution. According to McAllester's standard PAC-Bayes bound
\cite{mcallester1999some}, for any $\delta \in (0, 1)$, the expected target risk
under the posterior distribution $Q$ is bounded with probability at least
$1-\delta$ by:
\begin{equation}
\label{eq:pac_bayes_bound}
\begin{aligned}
\mathbb{E}_{\Lambda \sim Q} [R(\theta(\Lambda))] &\le \mathbb{E}_{\Lambda \sim
Q} [\hat{R}(\theta(\Lambda); X)] \\
&+ \sqrt{\frac{\text{KL}(Q(\Lambda) \| P(\Lambda)) + \log \frac{2\sqrt{N}}{\delta}}{2N}}
\end{aligned}
\end{equation}
While McAllester's bound in Eq.~\ref{eq:pac_bayes_bound} features a non-linear
square root over the complexity term, we can alternatively employ Alquier's
linear PAC-Bayes bound \cite{alquier2016properties}. For any temperature
parameter $\beta > 0$, the expected target risk under $Q$ is bounded with
probability at least $1-\delta$ by:
\begin{equation}
\label{eq:alquier_bound}
\begin{aligned}
\mathbb{E}_{\Lambda \sim Q} [R(\theta(\Lambda))] &\le \mathbb{E}_{\Lambda \sim
Q} [\hat{R}(\theta(\Lambda); X)] \\
&+ \frac{\text{KL}(Q(\Lambda) \| P(\Lambda)) + \ln(1/\delta)}{\beta N} + \frac{\beta}{8}
\end{aligned}
\end{equation}
This formulation is highly advantageous for our optimization scheme. Because
Eq.~\ref{eq:alquier_bound} is strictly linear with respect to the KL divergence,
minimizing Alquier's bound with respect to $Q$ is mathematically equivalent to
minimizing the sum of the empirical risk and a linear KL complexity penalty.
This provides a direct, mathematically rigorous justification for our final
linear-in-KL optimization objective.

\begin{remark}[Mathematical Necessity of Class-Capacity Normalization]
McAllester's standard PAC-Bayes bound in Eq.~\ref{eq:pac_bayes_bound} holds
strictly for empirical risks bounded in the interval $[0, 1]$. While standard
Shannon entropy $\mathcal{H}_k \in [0, \log C_k]$ is unbounded as the number of
classes $C_k$ grows, our Class-Capacity Normalization (CCN) scales the entropy
by the maximum possible value $\log C_k$. This forces our test-time adaptative
loss $\tilde{\mathcal{H}}_k = \frac{\mathcal{H}_k}{\log C_k}$ to reside strictly
within $[0, 1]$. Thus, CCN is not merely an empirical calibration heuristic, but
a fundamental mathematical requirement that guarantees the formal validity of
the PAC-Bayes generalization bound over test-time entropy minimization.
\end{remark}

Assuming independence across tasks, the prior and posterior factorize as
$P(\Lambda) = \prod_{k=1}^K P(\lambda_k)$ and $Q(\Lambda) = \prod_{k=1}^K
Q(\lambda_k)$, so that the complexity penalty decomposes into:
\begin{equation}
\text{KL}(Q(\Lambda) \| P(\Lambda)) = \sum_{k=1}^K \text{KL}(Q(\lambda_k) \| P(\lambda_k))
\label{eq:kl_decomposition}
\end{equation}
We define the prior distribution for each task coefficient vector $P(\lambda_k)
= \mathcal{N}(\mu_0, \Sigma_{\ell})$ as a multivariate Gaussian centered at a
uniform task-arithmetic prior mean $\mu_0 = \frac{1}{K} \mathbf{1}_L$. We
parameterize the posterior distribution $Q(\lambda_k) = \mathcal{N}(\lambda_k^*,
\sigma_q^2 I)$ as an isotropic Gaussian centered at the optimized coefficients
$\lambda_k^* \in \mathbb{R}^L$ with a small fixed variance $\sigma_q^2$.

We address the deep mathematical and statistical foundations of this formulation
in the Appendix. Specifically, in Appendix~\ref{sec:appendix_remarks} we
resolve: (1) \textbf{The Truncated Gaussian Paradox} (avoiding the KL explosion
through narrow posterior variance and Lipschitz continuity); (2) \textbf{The
Surrogate-to-Target Risk Gap} (unsupervised entropy-minimization limits); and
(3) \textbf{Boundary Truncation Bias} (proving that partition function
truncation gradients are negligible under projected clamping, validating our
simple quadratic objective).

The KL divergence between two multivariate Gaussian distributions in $\mathbb{R}^L$ is:
\begin{equation}
\label{eq:kl_gaussian}
\begin{aligned}
\text{KL}(Q(\lambda_k) &\| P(\lambda_k)) = \frac{1}{2} \bigg[
\text{Tr}\left(\Sigma_{\ell}^{-1} (\sigma_q^2 I)\right) \\
&+ (\lambda_k^* - \mu_0)^T \Sigma_{\ell}^{-1} (\lambda_k^* - \mu_0) \\
&- L + \ln \frac{\det \Sigma_{\ell}}{\det(\sigma_q^2 I)} \bigg]
\end{aligned}
\end{equation}
Since $\mu_0$, $\Sigma_{\ell}$, $L$, and $\sigma_q^2$ are constants with respect
to the optimization parameters $\lambda_k^*$, minimizing Alquier's linear
PAC-Bayes complexity bound with respect to $\Lambda^* = \{\lambda_k^*\}$ is
mathematically equivalent to minimizing the quadratic penalty:
\begin{equation}
\mathcal{L}_{\text{GP}}(\Lambda^*) = \frac{1}{2} \sum_{k=1}^K (\lambda_k^* -
\mu_0)^T \Sigma_{\ell}^{-1} (\lambda_k^* - \mu_0)
\label{eq:gp_penalty}
\end{equation}
This establishes a mathematically rigorous derivation for quadratic
precision-matrix regularization directly from PAC-Bayes theory!

\begin{remark}[Randomized-to-Deterministic Discrepancy]
We explicitly acknowledge a subtle theoretical-to-empirical discrepancy standard
in PAC-Bayes applications. The formal generalization bounds in
Eq.~\ref{eq:pac_bayes_bound} and Eq.~\ref{eq:alquier_bound} bound the expected
target risk under a randomized posterior classifier $\Lambda \sim Q$. In
practice, however, model merging is evaluated deterministically using the
optimized mean coefficients $\Lambda^*$. We justify evaluating the single
deterministic model $\theta(\Lambda^*)$ by utilizing the narrow posterior
variance $\sigma_q^2 \to 0$ and the Lipschitz continuity of the loss landscape.
As $\sigma_q^2$ vanishes, the posterior distribution $Q(\Lambda)$ collapses to a
Dirac delta centered at $\Lambda^*$, which aligns the expected risk under $Q$
with the deterministic risk at $\Lambda^*$. While this deterministic evaluation
is a highly practical and standard approximation, the strict theoretical
generalization bounds formally apply to the randomized ensemble. Specifically,
the deterministic model acts as a point estimator of the ensemble behavior;
while our empirical evaluations demonstrate that this approximation yields
exceptional performance, we caution that the strict mathematical guarantees are
technically softened when moving from the stochastic ensemble to a single
deterministic parameter state.
\end{remark}

\subsection{The Spatial Gaussian Process Prior}

To model the physical architecture of neural networks, we place a Gaussian
Process (GP) prior over the layer-wise coefficients as a function of depth. Let
$z_l = l/L \in [0, 1]$ represent the normalized layer coordinate of layer $l$.
The spatial correlation between layers is modeled using a Squared Exponential
(RBF) kernel, defining the GP covariance matrix $\Sigma_{\ell} \in \mathbb{R}^{L
\times L}$:
\begin{equation}
[\Sigma_{\ell}]_{l, l'} = \sigma_p^2 \exp\left( - \frac{(z_l -
z_{l'})^2}{2\ell^2} \right) + \sigma_n^2 \delta_{l, l'}
\label{eq:gp_kernel}
\end{equation}
where $\sigma_p^2$ is the prior signal variance, $\ell$ is the spatial
lengthscale, and $\sigma_n^2$ is a small diagonal noise term (jitter) that
guarantees numerical invertibility of $\Sigma_{\ell}$.

\subsubsection{Theoretical Properties of the Precision Matrix $\Sigma_{\ell}^{-1}$}

The precision matrix $\Sigma_{\ell}^{-1}$ is a positive-definite, symmetric
operator. Its quadratic form in Eq. \ref{eq:gp_penalty} reveals two elegant and
unified physical constraints:
\begin{enumerate}
    \item \textbf{Proximity Penalty (Diagonal Entries):} The diagonal elements
    $[\Sigma_{\ell}^{-1}]_{l, l} > 0$ penalize large deviations of any layer's
    coefficient from the task-arithmetic prior mean $\mu_{0}$. This acts as a
    localized weight decay, bounding parameter drift at test-time.
    \item \textbf{Spatial Smoothness (Off-Diagonal Entries):} Because
    neighboring layers are highly correlated in the prior ($[\Sigma_{\ell}]_{l,
    l+1} > 0$), the adjacent off-diagonal elements $[\Sigma_{\ell}^{-1}]_{l,
    l+1}$ are negative. While a general Squared Exponential (RBF) prior results
    in a dense precision matrix that couples all layers, its adjacent-layer
    terms dominate. To gain analytical intuition, we can approximate the
    quadratic form using its tridiagonal components, or equivalently, obtain
    this decomposition \emph{exactly} by placing an Ornstein-Uhlenbeck prior (as
    discussed in \cref{rem:rbf_vs_ou} below). Under this tridiagonal structure, the
    quadratic product expands to:
    \begin{equation}
    \label{eq:precision_expansion}
    \begin{aligned}
    (\lambda_k - \mu_0)^T &\Sigma_{\ell}^{-1} (\lambda_k - \mu_0) \approx
    \sum_{l=1}^L a_l (\lambda_{l, k} - \mu_{0, l})^2 \\
    &+ \sum_{l=1}^{L-1} b_l (\lambda_{l, k} - \lambda_{l+1, k})^2
    \end{aligned}
    \end{equation}
    where $a_l, b_l > 0$. The second term acts as a finite-difference Laplacian
    smoother, directly penalizing high-frequency spatial volatility across
    adjacent layers and resolving the Overfitting-Optimizer Paradox.
\end{enumerate}

\begin{remark}[RBF vs. Ornstein-Uhlenbeck Precision Matrices]
\label{rem:rbf_vs_ou}
We highlight an elegant mathematical distinction between GP covariance kernels.
The strictly tridiagonal adjacent-layer coupling in
Eq.~\ref{eq:precision_expansion} is mathematically exact under the
\textbf{Ornstein-Uhlenbeck (OU) process}, which uses the exponential kernel:
\begin{equation}
[\Sigma_{\text{OU}}]_{l, l'} = \sigma_p^2 \exp\left( - \frac{|z_l -
z_{l'}|}{\ell} \right) + \sigma_n^2 \delta_{l, l'}
\label{eq:ou_kernel}
\end{equation}
Since the OU process is a first-order Markov process, its precision matrix
$\Sigma_{\text{OU}}^{-1}$ is strictly tridiagonal, coupling only adjacent
layers. Under our Squared Exponential (RBF) kernel (Eq.~\ref{eq:gp_kernel}), the
precision matrix $\Sigma_{\ell}^{-1}$ is dense (coupling all layers), but its
off-diagonal entries decay exponentially away from the diagonal. This dense
precision matrix acts as a higher-order continuous spatial smoother, penalizing
both adjacent and multi-hop layer differences. Both kernels successfully
stabilize test-time model merging, but the OU kernel provides a computationally
cheaper tridiagonal structure. Specifically, while computing the dense inverse
of the RBF kernel scales cubically as $O(L^3)$, the OU precision matrix
$\Sigma_{\text{OU}}^{-1}$ has a closed-form analytical tridiagonal expression
that can be computed in linear time $O(L)$ without any matrix inversion,
offering perfect scalability for ultra-deep architectures (e.g., models with
hundreds or thousands of layers).
\end{remark}

\subsubsection{Landscape-Smoothing and Convexifying Geometric Effect}

We highlight an important optimization advantage of our quadratic
precision-matrix penalty. The unsupervised test-time surrogate loss landscape
(such as prediction entropy) is highly non-convex and cluttered with numerous
local minima. In such environments, unconstrained first-order optimizers
(e.g., Adam) are easily trapped in poor local basins, or they oscillate wildly
when fitting transductive noise.

Adding our GP-prior penalty introduces a strongly convex quadratic term
$\frac{\alpha}{2} (\lambda^* - \mu_0)^T \Sigma_{\ell}^{-1} (\lambda^* - \mu_0)$
to the overall loss objective. Because the precision matrix $\Sigma_{\ell}^{-1}$
is positive-definite, this quadratic term behaves geometrically as a
multi-dimensional paraboloid centered at the prior mean. This convex quadratic
term effectively ``convexifies'' and smooths out the local, high-frequency
non-convex oscillations of the entropy landscape. It acts as a powerful
pre-conditioner that regularizes the gradients, allowing first-order
optimization to escape shallow, noisy local minima and converge rapidly to the
optimal global trajectory. This geometric smoothing effect directly explains the
outstanding empirical convergence speed observed in our experiments (reported
in Section 4), where GP-BayesMerge achieves near-peak performance in fewer than
50 gradient steps.

\subsubsection{Continuous Interpolation Limits}

The GP lengthscale $\ell$ provides a continuous, smooth interpolation across
three distinct merging regimes:

\textbf{Independent Weight Decay ($\ell \to 0$):} As $\ell \to 0$, neighboring
layers decouple, and $\Sigma_{\ell} \to (\sigma_p^2 + \sigma_n^2) I$. The
regularization reduces to a standard decoupled $L_2$ proximity penalty
(independent layer weight decay):
\begin{equation}
\lim_{\ell \to 0} \mathcal{L}_{\text{GP}}(\Lambda^*) = \frac{1}{2(\sigma_p^2 +
\sigma_n^2)} \sum_{k=1}^K \|\lambda_k^* - \mu_0\|_2^2
\end{equation}

\textbf{Flat Spatial Averaging ($\ell \to \infty$):} As $\ell \to \infty$, all
layers become perfectly correlated, forcing $\lambda_{1, k} = \lambda_{2, k} =
\dots = \lambda_{L, k}$. The precision matrix heavily penalizes any deviation
from a flat spatial average across layers, forcing constant coefficients per
task.

\textbf{Spatially Correlated Heterogeneity (Intermediate $\ell$):} For
intermediate values (such as $\ell \in [0.1, 0.3]$), GP-BayesMerge allows
coefficients to vary smoothly across layers to capture localized architectural
trends, while pruning high-frequency transductive noise.

\subsubsection{Lengthscale Scaling and Architectural Stability as $L \to \infty$}
\label{subsec:lengthscale_scaling}

A crucial consideration is the behavior of the GP prior as the neural network
depth $L \to \infty$ (e.g., in extremely deep foundation models). Since the
covariance matrix $\Sigma_{\ell}$ is defined over normalized coordinates $z_l =
l/L \in [0, 1]$, the distance between adjacent layers scales as $1/L$. If the
lengthscale $\ell$ is held constant as $L \to \infty$, the adjacent-layer
correlation approaches $\exp(0) = 1$. In extremely deep models, this causes
neighboring layers to become nearly perfectly correlated, leading to severe
multicollinearity and ill-conditioning of the covariance matrix $\Sigma_{\ell}$.

To prevent matrix ill-conditioning and preserve structural stability, the
lengthscale should scale inversely with the network depth $L$. Specifically, to
maintain a stable physical correlation distance (e.g., correlating layers within
a constant physical block size of $B_{\text{phys}}$ layers), we propose the
scaling rule:
\begin{equation}
\ell = \frac{B_{\text{phys}}}{L}
\label{eq:lengthscale_scaling}
\end{equation}
By setting $\ell \propto 1/L$, the normalized correlation distance $\frac{\lvert
z_l - z_{l'} \rvert}{\ell} = \frac{\lvert l - l' \rvert}{B_{\text{phys}}}$
becomes independent of the total depth $L$, successfully stabilizing the
physical correlation distance across architectures and preventing numerical
instabilities as networks grow deeper.

\subsubsection{Non-Stationary Block-Wise GP Prior Boundaries}
\label{subsubsec:block_wise_gp}

While a continuous, stationary Gaussian Process prior over normalized network
depth is highly effective, modern deep architectures (such as Vision
Transformers or ResNets) are fundamentally non-stationary. They are partitioned
into distinct functional stages or block types (e.g., self-attention blocks
versus multilayer perceptrons, or residual blocks operating at different feature
resolutions). Across these functional boundaries, we expect sharp architectural
transitions where the optimal merging coefficients may change abruptly. Force-fitting
a stationary GP prior over the entire network depth can cause over-smoothing
across these boundaries, dampening localized transitions.

To handle these highly heterogeneous, multi-stage architectures, we propose a
\emph{non-stationary block-wise GP prior}. We partition the $L$ layers of the
network into $M$ distinct functional blocks. Let $b(l) \in \{1, \dots, M\}$
denote the block assignment of layer $l$. We define the block-wise covariance
matrix $\Sigma_{\text{block}} \in \mathbb{R}^{L \times L}$ by introducing a
decoupling factor $\rho \in [0, 1]$ that scales the covariance between layers
residing in different functional blocks:
\begin{equation}
[\Sigma_{\text{block}}]_{l, l'} = 
\begin{cases}
[\Sigma_{\ell}]_{l, l'} & \text{if } b(l) = b(l') \\
\rho [\Sigma_{\ell}]_{l, l'} & \text{if } b(l) \neq b(l')
\end{cases}
\label{eq:block_covariance}
\end{equation}
Here, $\Sigma_{\ell}$ is the base stationary covariance matrix (Eq.~\ref{eq:gp_kernel}).
Within each functional block, layers remain highly correlated (governed by
lengthscale $\ell$), preserving spatial smoothness. Across block boundaries,
however, the correlation is scaled down by the decoupling factor $\rho$, allowing
coefficients to decouple and adapt to sharp, block-specific transitions on-the-fly.
This formulation generalizes our stationary GP model: when $\rho = 1.0$, it
collapses back to the stationary GP prior, while setting $\rho = 0.0$ yields a
fully block-diagonal covariance where different stages are optimized in complete
isolation. As shown in our empirical analyses (Section 4 and Appendix C.1), this
non-stationary block-wise prior successfully prevents over-smoothing across
functionally distinct stages, achieving superior peak performance on highly
heterogeneous target tasks.

\subsection{Theoretical Generalization: Kronecker Multi-Task GP Priors}
\label{subsec:multitask_gp}

Our formulation in Eq.~\ref{eq:gp_penalty} assumes that the prior and posterior
distributions over the merging coefficients factorize independently across the
$K$ tasks, i.e., $P(\Lambda) = \prod_{k=1}^K P(\lambda_k)$ and $Q(\Lambda) =
\prod_{k=1}^K Q(\lambda_k)$. While computationally lightweight and highly
effective in practice, this independent task assumption ignores representational
conflicts, sign interferences, and parameter trade-offs that often occur when
merging highly distinct task experts. In actual model merging, the optimal
coefficient profile for Task $A$ is inherently coupled to that of Task $B$ due
to shared network parameters.

To resolve this limitation from first principles, we present a mathematically
rigorous generalization of GP-BayesMerge to a joint, multi-task Gaussian Process
prior. Let $\lambda = \text{vec}(\Lambda) = [ \lambda_1^T, \dots, \lambda_K^T ]^T \in \mathbb{R}^{LK}$
be the vectorized joint coefficient vector obtained by concatenating the layer-wise coefficients of all tasks.
We define the corresponding optimized vectorized joint coefficient vector as $\lambda^* = \text{vec}(\Lambda^*) = [ (\lambda_1^*)^T, \dots, (\lambda_K^*)^T ]^T \in \mathbb{R}^{LK}$. We place a joint multivariate Gaussian prior over the entire
coefficient space:
\begin{equation}
P( \lambda ) = \mathcal{N}\left( \mu_{\text{joint}}, \Sigma_{\text{joint}} \right)
\end{equation}
where $\mu_{\text{joint}} = \frac{1}{K} \mathbf{1}_{LK}$ is the joint uniform
prior mean, and the joint covariance matrix $\Sigma_{\text{joint}} \in
\mathbb{R}^{LK \times LK}$ is modeled as the Kronecker product of a task
correlation matrix $B \in \mathbb{R}^{K \times K}$ and our spatial depth
covariance matrix $\Sigma_{\ell} \in \mathbb{R}^{L \times L}$:
\begin{equation}
\Sigma_{\text{joint}} = B \otimes \Sigma_{\ell}
\end{equation}
Here, the task-similarity matrix $B \in \mathbb{R}^{K \times K}$ is a symmetric,
positive-definite matrix that encodes the known semantic or representational
relationships between tasks (e.g., computed offline via activation CKA
similarities or Fisher Information overlaps between task experts). By utilizing
the properties of the Kronecker product, the joint precision matrix factors
analytically as:
\begin{equation}
\Sigma_{\text{joint}}^{-1} = B^{-1} \otimes \Sigma_{\ell}^{-1}
\label{eq:kronecker_precision}
\end{equation}
This factorization is highly powerful: it completely bypasses the $O(L^3 K^3)$
cubic inversion of the joint covariance matrix, reducing the computation to
separate inversions of $B$ ($O(K^3)$) and $\Sigma_{\ell}$ ($O(L^3)$).

The corresponding multi-task joint GP complexity penalty simplifies to the quadratic form:
\begin{equation}
\label{eq:multitask_gp_penalty}
\begin{aligned}
\mathcal{L}_{\text{MT-GP}}(\Lambda^*) &= \frac{1}{2} (\lambda^* - \mu_{\text{joint}})^T \\
&\quad \times \left( B^{-1} \otimes \Sigma_{\ell}^{-1} \right) (\lambda^* - \mu_{\text{joint}})
\end{aligned}
\end{equation}
This joint multi-task formulation penalizes high-frequency noise not only across
adjacent layers (via $\Sigma_{\ell}^{-1}$) but also across conflicting task
directions (via $B^{-1}$), providing a unified and mathematically elegant
framework for structured, multi-task test-time weight interpolation.

\begin{remark}[Relaxing the Offline Data Assumption for $B$ via Online Calibration]
While computing the task-similarity matrix $B$ offline using the experts'
original training or validation datasets provides a robust prior, it assumes
access to original training data which might be prohibited due to privacy or
bandwidth constraints in real-world test-time edge deployments. To relax this
offline assumption and ensure a fully training-free and data-free test-time
adaptation, we can estimate $B$ dynamically and online. Specifically, as the
small target calibration batches $\{X_k\}$ arrive at test-time, we feed these
unlabeled samples through each of the pre-trained expert models to extract their
intermediate representations. We then compute the pairwise Centered Kernel
Alignment (CKA) similarity between task experts' activations \emph{on-the-fly}
using only the test-time calibration samples. This online CKA estimation is
fully data-free, adds negligible computational overhead (as it scales
quadratically with $K \ll D$), and dynamically adapts the task correlation prior
to the actual target distribution encountered at test-time, thereby resolving
representational conflicts under covariate shifts without requiring any offline
dataset access. Crucially, under non-stationary streaming scenarios with temporal domain drift, $B_{\text{online}}$ can be sequentially updated on-the-fly using a sequential Bayesian filter or a sliding temporal window over activation CKA features (e.g., $B_{t} = (1 - \rho) B_{t-1} + \rho B_{\text{batch}, t}$, where $\rho \in (0, 1)$ is a temporal decay parameter). This ensures that task correlations adapt dynamically and smoothly to non-stationary environments, preventing representational lag on dynamically evolving tasks without introducing any computational latency.

Furthermore, we address the stability of the online-estimated matrix $B$ under
extremely low-sample regimes (e.g., $N \in \{4, 8\}$). While small sample sizes
can introduce transductive noise into individual activation similarities, our
Kronecker product factorization $\Sigma_{\text{joint}}^{-1} = B^{-1} \otimes
\Sigma_{\ell}^{-1}$ provides an exceptionally robust shielding mechanism.
Specifically, the spatial GP precision matrix $\Sigma_{\ell}^{-1}$ acts as the
dominant global smoother across network depth. Even if a noisy $B$ introduces
minor distortions in task-wise coupling, high-frequency layer-specific
oscillations are still heavily penalized by $\Sigma_{\ell}^{-1}$, preventing
optimization divergence. Consequently, MT-GP-BayesMerge remains highly stable
under low-sample regimes, degrading gracefully without losing its competitive
edge.

To guarantee absolute numerical stability during the matrix inversion $B^{-1}$
in low-sample regimes where $B$ could potentially become close-to-singular, we
apply a small, positive diagonal jitter (shrinkage) to the estimated correlation
matrix. Specifically, we define a stabilized correlation matrix
$B_{\text{stable}} = (1 - \epsilon) B_{\text{online}} + \epsilon I$, where $I
\in \mathbb{R}^{K \times K}$ is the identity matrix and $\epsilon = 10^{-4}$
represents a small regularization multiplier. This simple, training-free
shrinkage operation places a rigorous lower bound on the eigenvalues of the
joint task covariance matrix, mathematically guaranteeing that the task
correlation matrix remains invertible and exceptionally well-conditioned under
any arbitrary test-time calibration batch, thereby preventing any numerical
overflow or gradient explosion.
\end{remark}

\begin{remark}[Theoretical Justification of Activation CKA over Parameter-Space Sim]
\label{rem:cka_justification}
We provide a formal justification for why online activation-based Centered Kernel Alignment (CKA) similarity \cite{kornblith2019similarity} is a highly sufficient and superior proxy for task-correlation $B$ compared to direct weight-space distance or parameter-level gradient overlaps (like Fisher Information overlaps):
\begin{enumerate}
    \item \textbf{Invariance to Model Symmetries:} Neural networks exhibit massive representation symmetries, such as node permutation and orthogonal coordinate rotations in weight-space. Two expert models can possess identical input-output mappings (and thus high functional compatibility) while residing in vastly different regions of parameter-space. Direct parameter-space comparisons (like $L_2$ weight distances) fail to capture these functional symmetries, falsely penalizing compatible models. Activation CKA, being invariant to orthogonal transformations and scaling of representations, bypasses these weight-space coordinate dependencies to measure functional representational alignment directly.
    \item \textbf{High-Dimensional Computational Tractability:} Calculating parameter-level Fisher Information overlaps requires computing and storing dense Fisher matrices or high-dimensional gradient vectors of size $O(D \times D)$ or $O(D)$ where $D$ is in millions or billions of parameters. This introduces catastrophic computational and memory overheads, making it physically impossible to execute on edge devices during real-time test-time adaptation. In contrast, online activation CKA is computed over the intermediate activations of size $O(N \times d_{\text{feat}})$ where $N \le 8$ is the calibration batch size and $d_{\text{feat}} \ll D$ is the hidden feature dimension. CKA is thus highly lightweight and computationally instantaneous.
    \item \textbf{Direct Capture of Representational Interference:} Model merging ultimately aims to combine experts such that they do not destructively interfere in their intermediate feature representations. Activation similarity directly measures this intermediate interference: if two experts yield highly aligned activation subspaces on target inputs, their features can be smoothly interpolated without representational collapse, ensuring that activation CKA is a direct, robust, and physically grounded proxy for merging compatibility.
\end{enumerate}
\end{remark}

\subsection{Unsupervised Test-Time Adaptation and the Surrogate-to-Target Risk Gap}
\label{subsec:surrogate_gap}

We recognize a fundamental theoretical limitation in unsupervised test-time
adaptation (TTA), which we term the \emph{surrogate-to-target risk gap}.
McAllester's and Alquier's PAC-Bayes bounds (Eq.~\ref{eq:pac_bayes_bound} and
Eq.~\ref{eq:alquier_bound}) require that the empirical risk $\hat{R}$ and target
risk $R$ are evaluated under the same loss function. Since true labels are
unavailable at test-time, the model minimizes Class-Capacity Normalized (CCN)
prediction entropy $\tilde{\mathcal{H}}_k$ as an unsupervised surrogate.

However, bounding expected entropy on the target distribution does \emph{not}
mathematically bound the expected classification error (true risk); under severe
covariate or domain shifts, a network can confidently predict incorrect classes,
resulting in low entropy but $100\%$ classification error. Thus, the theoretical
guarantees strictly apply only to the entropy surrogate, not the target
classification accuracy itself.

To mathematically bridge the surrogate-to-target risk gap, we establish two
formal theoretical conditions under which minimizing the prediction entropy
surrogate is guaranteed to reduce classification risk. These assumptions are
grounded in semi-supervised learning and domain adaptation theory:
\begin{enumerate}
    \item \textbf{The Margin-Preserving Support Assumption:} We assume that the
    target domain's features lie in regions where the decision boundaries of the
    source classifiers are well-defined. That is, the true optimal
    classification boundary resides in low-density valleys of the target feature
    distribution, which is the foundational support of the Cluster Assumption in
    semi-supervised learning. Under this assumption, pushing predictions away
    from boundaries (reducing prediction entropy) pushes target features into
    high-density class regions, aligning entropy reduction with error reduction.
    \item \textbf{Classifier Calibration Assumption:} We assume that the task
    expert classifiers are well-calibrated on the target domain. Specifically,
    let $p(y=c | x)$ be the true posterior probability and $\hat{p}(y=c | x)$ be
    the predicted probability. Calibration requires that as $\max_c
    \hat{p}(y=c|x) \to 1$, the empirical probability of correctness also
    approaches $1$. If the expert heads are calibrated, high-confidence (low
    entropy) predictions are statistically guaranteed to be correct,
    mathematically bounding classification error.
\end{enumerate}
To mathematically bridge this gap, we establish a formal theoretical result that
bounds the expected true classification risk by the expected prediction entropy
under our two formal assumptions.

\begin{theorem}[Surrogate-to-Target Risk Bound]
\label{thm:surrogate_to_target_bound}
Let $f(x) \in \Delta^{C-1}$ be the model's predicted probability vector over $C$
classes for target input $x$, let $g(x) = \max_c f_c(x)$ be the prediction
confidence, and let $\hat{y}(x) = \arg\max_c f_c(x)$ be the predicted label.
Suppose the target domain distribution satisfies:
\begin{enumerate}
    \item \textbf{Margin-Preserving Support (with margin parameter $\gamma \in
    [\frac{1}{2}, 1]$):} The confidence satisfies $g(x) \ge \gamma$ almost
    surely over the target distribution.
    \item \textbf{Expected Calibration Error Bound:} The pointwise calibration
    discrepancy $\epsilon(x) = \lvert P(y \neq \hat{y}(x) \mid x) - (1 - g(x))
    \rvert$ satisfies $\mathbb{E}_{x \sim \mathcal{D}_k}[\epsilon(x)] \le
    \mathcal{E}_{\text{cal}}$.
\end{enumerate}
Then, the true classification risk $R(f) = \mathbb{E}_{x \sim
\mathcal{D}_k}[\mathbb{1}(\hat{y}(x) \neq y)]$ is upper-bounded by the expected
prediction entropy as:
\begin{equation}
R(f) \le \frac{\ln C}{2\gamma} \mathbb{E}_{x \sim
\mathcal{D}_k}[\tilde{\mathcal{H}}(f(x))] + \mathcal{E}_{\text{cal}}
\end{equation}
where $\tilde{\mathcal{H}}(f(x)) = \frac{\mathcal{H}(f(x))}{\ln C}$ is the normalized entropy.\footnote{We discuss the practical limitations of this theorem's assumptions under severe out-of-distribution shifts and propose dynamic, data-driven relaxations for the margin parameter $\gamma$ in the \emph{Discussion on the Practical and Theoretical Limits of Unsupervised TTA} paragraph below.}
\end{theorem}

\begin{proof}
Let $f_c(x)$ be the predicted probability of class $c$. The entropy is
$\mathcal{H}(f(x)) = -\sum_{c=1}^C f_c(x) \ln f_c(x)$. Using the inequality
$-\ln u \ge 1-u$ for all $u > 0$, we have:
\begin{equation}
\mathcal{H}(f(x)) \ge \sum_{c=1}^C f_c(x)(1 - f_c(x)) = 1 - \sum_{c=1}^C f_c(x)^2
\end{equation}
Let $\hat{y} = \arg\max_c f_c(x)$ and $g(x) = f_{\hat{y}}(x)$. Since $f_c(x) \ge
0$ and $\sum_{c \neq \hat{y}} f_c(x) = 1 - g(x)$, we have by the Cauchy-Schwarz
inequality:
\begin{equation}
\sum_{c \neq \hat{y}} f_c(x)^2 \le \left(\sum_{c \neq \hat{y}} f_c(x)\right)^2 = (1 - g(x))^2
\end{equation}
Thus, the sum of squares is bounded as:
\begin{equation}
\begin{aligned}
\sum_{c=1}^C f_c(x)^2 &= g(x)^2 + \sum_{c \neq \hat{y}} f_c(x)^2 \\
&\le g(x)^2 + (1 - g(x))^2
\end{aligned}
\end{equation}
Expanding this gives $\sum_{c=1}^C f_c(x)^2 \le 2g(x)^2 - 2g(x) + 1$. Plugging
this into the entropy lower bound yields:
\begin{equation}
\mathcal{H}(f(x)) \ge 1 - (2g(x)^2 - 2g(x) + 1) = 2g(x)(1 - g(x))
\end{equation}
Under the Margin-Preserving Support assumption, the target inputs lie in regions
where $g(x) \ge \gamma \ge \frac{1}{2}$. Thus, $2g(x) \ge 2\gamma \ge 1$, which
implies:
\begin{equation}
1 - g(x) \le \frac{1}{2\gamma} \mathcal{H}(f(x)) = \frac{\ln C}{2\gamma} \tilde{\mathcal{H}}(f(x))
\end{equation}
Now, we decompose the expected true classification risk $R(f) = \mathbb{E}_x
[P(y \neq \hat{y}(x) \mid x)]$. Since $P(y \neq \hat{y}(x) \mid x) \le 1 - g(x)
+ \epsilon(x)$, taking expectations on both sides gives:
\begin{equation}
\begin{aligned}
R(f) &\le \mathbb{E}_{x \sim \mathcal{D}_k}[1 - g(x)] + \mathbb{E}_{x \sim
\mathcal{D}_k}[\epsilon(x)] \\
&\le \frac{\ln C}{2\gamma} \mathbb{E}_{x \sim
\mathcal{D}_k}[\tilde{\mathcal{H}}(f(x))] + \mathcal{E}_{\text{cal}}
\end{aligned}
\end{equation}
which completes the proof.
\end{proof}

\paragraph{Discussion on the Practical and Theoretical Limits of Unsupervised TTA.}
To maintain radical transparency, we must emphasize that while
Theorem~\ref{thm:surrogate_to_target_bound} provides a formal mathematical
bridge between prediction entropy and true classification error, its practical
guarantees depend heavily on the validity of its underlying conditions.
Specifically, the Expected Calibration Error (ECE) bound
$\mathcal{E}_{\text{cal}}$ must remain small. Under severe domain shifts or
catastrophic covariate shifts where the model is uncalibrated on the target
domain, $\mathcal{E}_{\text{cal}}$ can become arbitrarily large. In such
uncalibrated regimes, minimizing prediction entropy can lead to self-reinforcing
"confirmation bias" errors: the model confidently predicts incorrect classes,
resulting in low entropy but $100\%$ classification error.

Furthermore, the Margin-Preserving Support assumption requires that the prediction confidence satisfies $g(x) \ge \gamma \ge 0.5$ almost surely over the target distribution. While this condition holds when the target domain is relatively structured or close-to-distribution, severe domain shifts can collapse classifier confidence, resulting in $g(x) < 0.5$ (e.g., approaching $1/C$ for uniform predictions). In such low-confidence regimes, the Margin-Preserving Support condition is violated, and naive entropy minimization can easily amplify prediction errors. To relax this restriction, practitioners can dynamically validate or adaptively set the margin parameter $\gamma$ on the incoming test stream. For example, instead of keeping $\gamma$ fixed, $\gamma$ can be set dynamically as the $\beta_{\text{conf}}$-quantile of the model's confidences over the incoming batch: $\gamma_k = \text{quantile}(\{g(x) : x \in X_k\}, \beta_{\text{conf}})$. This dynamic scaling allows the margin to adapt automatically to severe OOD confidence collapses, restricting adaptation gradients to the relatively most confident target subset. Alternatively, practitioners can employ dynamic confidence thresholding (e.g., sourcing adaptation gradients exclusively from samples where $g(x) \ge \gamma_k$, while ignoring highly uncertain inputs) to ensure that decision boundaries are adjusted based only on well-separated, cluster-pure target regions. Consequently, practitioners must remain aware that the unsupervised PAC-Bayes generalization bounds strictly govern only the entropy surrogate, and their empirical translation to target accuracy is contingent on the calibration and margin-preserving properties of the underlying task expert models.

\subsection{Unified Multi-Task Optimization Objective}

At test-time, the merged model is adapted on a small calibration batch $X =
\{X_1, \dots, X_K\}$ consisting of $N$ unlabeled samples per task. We utilize
prediction entropy as our unsupervised empirical risk. To prevent "sacrificial
task bias" on complex datasets, we integrate the Class-Capacity Normalization
(CCN) and Scale-Normalized Entropy Weighting (SNEW) calibration schemes
\cite{trial2_submission1}.

Let $C_k$ be the number of classes in task $k$. The Class-Capacity Normalized entropy is:
\begin{equation}
\tilde{\mathcal{H}}_k(\theta(\Lambda); X_k) = \frac{\mathcal{H}_k(\theta(\Lambda); X_k)}{\log C_k}
\label{eq:ccn_entropy}
\end{equation}
To balance the gradients of tasks with widely different baseline entropy levels,
SNEW assigns task weights $w_k$ inversely proportional to their baseline
normalized entropy under uniform task arithmetic $\tilde{\mathcal{H}}_k^0$:
\begin{equation}
w_k = \frac{1}{\tilde{\mathcal{H}}_k^0}
\label{eq:snew_weight}
\end{equation}
Combining the calibrated entropy loss with our PAC-Bayes GP regularization, the
complete test-time adaptation objective optimized by GP-BayesMerge is:
\begin{equation}
\label{eq:total_objective}
\begin{aligned}
\mathcal{L}_{\text{total}}(\Lambda^*) &= \sum_{k=1}^K w_k \cdot
\tilde{\mathcal{H}}_k(\theta(\Lambda^*); X_k) \\
&+ \frac{\alpha}{2} \sum_{k=1}^K (\lambda_k^* - \mu_0)^T \Sigma_{\ell}^{-1} (\lambda_k^* - \mu_0)
\end{aligned}
\end{equation}
where $\alpha > 0$ is the regularization scaling hyperparameter. The
optimization is performed via Adam gradient descent for $T$ steps on raw
unclamped coefficients $\Lambda^*_{\text{raw}}$, where the physical merging
coefficients evaluated by the network are obtained by clamping $\Lambda^* =
\text{clamp}(\Lambda^*_{\text{raw}}, 0.0, 1.0)$ at each step to preserve valid Weight
interpolation, while the quadratic GP prior penalty is evaluated directly on the
unclamped $\Lambda^*_{\text{raw}}$ to prevent gradient saturation on the boundary.\footnote{Evaluating the prior penalty on unclamped coordinates ensures a continuous restoring force back to the prior mean, preventing permanently zeroed-out gradients. We provide a complete mathematical and optimization justification for this unclamped regularization choice in Appendix~\ref{sec:appendix_remarks}.}
